\section{Discussion}\label{sec:discussion}

\paragraph{What changed.} The first kernel and the new one execute the same instructions with the same frame
semantics and the same noise decision rules; they differ in the counter based generator that supplies the random
words, and in where the work happens: noise generation moved off the critical loop into its own
pipelined stage, frame hazards moved from hardware into the compiler's schedule, detector folding moved into a stage
that overlaps the next batch, and the cost of a launch is shared by thousands of batches instead of one. Each of these
is a standard technique; together they turn the interpreter of the first kernel into a pipeline whose rate the
compiled program alone determines.

\paragraph{Nodes without communication.} Our distributed statevector simulator on the same four cards spends much of
its time exchanging amplitudes through host memory, and more cards add capacity there rather than
speed~\cite{ali2026multifpga}. Sampling is the opposite case: with independent random streams per node the nodes do
not interact, and the four card rate is four times the one card rate (Table~\ref{tab:throughput}). The node model carries over
unchanged; what decides scaling is whether the workload needs data from other nodes. For workloads of this kind,
more cards in a host, or hosts with cards, add rate in proportion until the host can no longer collect the output.

\paragraph{Routability as the limit.} The synthesis estimates were never the obstacle: every version of the kernel that
was implemented reached one slot per cycle and an estimated clock above the 300~MHz target. What decided between a working and a
failing design was routing, and the cause was density rather than size: the failed designs used \NumPhiloxDspCardPct\,\%
of the DSP blocks available to user logic but packed them into the DSP columns of the regions they occupied, and
the conflicting nets of every failed run lay in the instance that shared SLR0 with the platform logic. A generator
built from adders and rotations, which spreads over ordinary logic, and the move away from SLR0 turned five hour runs
that failed into a two hour run that met timing. For HLS designs that replicate arithmetic units in
parallel, we found the placement of those units on the device to matter as much as their number.

\paragraph{Limits of the kernel.} The frame memory holds \NumQubitsMax{} qubits after renumbering, enough for the
rotated surface code up to distance 7 (\NumSlotsSevenQubits{} qubits) but not distance 9. The program memory holds
\NumProgMax{} slots, the measurement record \NumMeasMax{} measurements and the output stage \NumOutMax{} detectors and
observables. Larger codes need wider limits; the frame memory grows linearly with the qubit count, and the program
and record memories are on chip block RAM. The instruction set covers the gates, measurements, resets and noise channels
of the generated memory experiments (H, S, CX, CZ, SWAP, measurement and reset in three bases, measure and reset,
single qubit Pauli errors, single and two qubit depolarizing noise); other Clifford gates, correlated errors and
Pauli product measurements would need further opcodes. Like the first kernel, it uses the frame semantics that
follow from a Clifford circuit; it does not handle non Clifford gates.

\paragraph{Limits of the comparison.} The CPU baseline is Stim's own detector sampler with one process per hardware
thread, which is how Stim is used for large sampling campaigns; a multi threaded sampler built on Stim might do better
on the same CPU. The comparison concerns sampling only: compiling the circuit, deriving the detector error model and
decoding are done on the host in both cases and are not timed. Energy was not compared, because the host provides no
power measurement of the CPU that would match the board power readings of the cards. The repetition code has only five qubits, so
most slots of its program are NOPs that keep reuses of a qubit apart (Table~\ref{tab:programs}); its rate reflects
the schedule more than the circuit, and the larger circuits are the more representative ones.

\paragraph{Embedded use.} The kernel needs no host interaction while a launch runs and its output is a stream of
detector words. The same pipeline could therefore feed a decoder on the same device or on a device next to it,
replacing recorded test data with live samples at hardware rate, which would make the sampler a test source for the
FPGA decoders of Section~\ref{sec:related}. We have not built that combination.
